% Part: first-order-logic
% Chapter: lindstrom
% Section: introduction

\documentclass[../../../include/open-logic-section]{subfiles}

\begin{document}

\olfileid{mod}{lin}{int}
\section{પરિચય}

આ પ્રકરણમાં આપણે પ્રથમ-ક્રમ તર્કશાસ્ત્રનું લિન્ડસ્ટ્રોમ-લક્ષણન સાબિત
કરવાનું લક્ષ્ય રાખીએ છીએ: અમુક વધારાની મર્યાદાઓ હેઠળ સઘનતા અને અધોગામી
લેવેનહાઇમ--સ્કોલેમ પ્રમેયો જે મહત્તમ તર્કશાસ્ત્ર માટે સત્ય છે તે
પ્રથમ-ક્રમ તર્કશાસ્ત્ર છે
(\olref[fol][com][com]{thm:compactness} અને
\olref[fol][com][dls]{thm:downward-ls}). પહેલાં આપણે આ પ્રમેય જેમને લાગુ
પડે છે તે તર્કશાસ્ત્રોના સામાન્ય વર્ગનું વધુ વ્યાપક લક્ષણન જોઈએ. આપણે
પોતાને \emph{સંબંધાત્મક} ભાષાઓ પૂરતા મર્યાદિત રાખીશું, એટલે કે એવી
ભાષાઓ જેમાં માત્ર !!{predicate}s અને વ્યક્તિ-અચળો હોય, પણ કોઈ
!!{function}s ન હોય.

\end{document}
